% !PRINTED 390
\weberpage{390}

Um dies nachzuweisen, genügt es zu zeigen, dass durch die
Permutationen dieser Gruppe irgend eine Wurzel, etwa [1\ 3], in
jede andere (mit Ausnahme von [1\ 2]) übergehen kann. Nun
geht aber [1,\,3] % [sic] comma printed here, unlike the bracket pairs elsewhere on this page which are space-separated without a comma
durch Permutationen der Gruppe $S$, durch die
1, 2 ungeändert bleiben, in [1\ 4], \dots, [1\ 8] über; ebenso [2\ 3] in
[2\ 4], \dots, [2\ 8]. Es bleibt also noch zu zeigen, dass man [1\ 3]
auch in [2\ 3] und in [4\ 5] überführen kann. Dies zeigt aber
der Anblick der drei folgenden vollständigen Siebener-Systeme:

\[
\begin{aligned}
&[1\ 2]\ [1\ 3]\ [1\ 4]\ [1\ 5]\ [1\ 6]\ [1\ 7]\ [1\ 8]\\
&[1\ 2]\ [2\ 3]\ [2\ 4]\ [2\ 5]\ [2\ 6]\ [2\ 7]\ [2\ 8]\\
&[1\ 2]\ [4\ 5]\ [3\ 4]\ [3\ 5]\ [1\ 6]\ [1\ 7]\ [1\ 8].
\end{aligned}
\]

Die Gruppe $P$ ist hiernach zweifach transitiv. Sie kann
aber nicht mehr als zweifach transitiv sein. Denn lässt man
zwei Wurzeln ungeändert, so kann eine mit diesen beiden syzygetische
Wurzel nicht in eine azygetische übergeführt werden.
Durch Adjunction von zwei Wurzeln wird die Doppeltangentengleichung
reducibel. Es löst sich ein Factor $10^{\text{ten}}$ Grades ab,
dessen Wurzeln mit den beiden gegebenen einen \gesperrt{Steiner}'schen
Complex bilden (die Function $X$ im Satze 1., \S.\,101).

Unter den Divisoren der Gruppe $P$ ist besonders die
Gruppe $S$ vom Index 36 bemerkenswerth. Diese Gruppe ist noch
transitiv, weil durch sie [1\ 2] in jede beliebige Wurzel $[\varkappa\ \lambda]$
übergeführt werden kann. Sie ist aber nur noch einfach transitiv,
weil bei festgehaltenem [1\ 2] die Wurzel [1\ 3] nicht mehr
in [4\ 5] übergehen kann.

Durch Adjunction \gesperrt{einer} Wurzel einer Gleichung $36^{\text{sten}}$ Grades
wird also die Gruppe der Doppeltangentengleichung auf die
Gruppe einer allgemeinen Gleichung $8^{\text{ten}}$ Grades reducirt.

Setzt man z.\,B.

\[
v_1 = [1\ 2]\ [1\ 3]\ [1\ 4]\ [1\ 5]\ [1\ 6]\ [1\ 7]\ [1\ 8],
\]

und definirt entsprechend $v_2, v_3, \dots, v_8$, so ist jede symmetrische
Function dieser acht Grössen, z.\,B.

\[
u = v_1 + v_2 + v_3 + v_4 + v_5 + v_6 + v_7 + v_8,
\]

Wurzel einer Gleichung $36^{\text{sten}}$ Grades. Adjungirt man dem
Problem die Grösse $u$, so werden die $v_1, v_2, \dots, v_8$ die Wurzeln
einer Gleichung $8^{\text{ten}}$ Grades, die keinen Affect hat.
